% =============================================================================
% CONCLUSION - CITA Paper in ALKALI Format
% =============================================================================

\vspace{-2mm}
\section{Discussion and Implications}
\label{sec:discussion}

This paper introduces \textbf{CITA (Contrastive Instruction-Tuned Alignment)}, a framework enabling instruction-aware alignment in large language models. Unlike DPO~\cite{rafailov2023direct} and RLHF~\cite{ouyang2022training} which embed static alignment into weights, CITA allows models to dynamically adapt behavior based on natural language instructions at inference time.

\vspace{-2mm}
\subsection{Key Findings}
\label{sec:key_findings}

\begin{enumerate}[leftmargin=*,itemsep=2pt]
    \item \textbf{CITA achieves highest instruction-alignment efficiency}: Using average normalized radar radius (\cref{eq:efficiency}), CITA achieves \textbf{86.7\%} efficiency---outperforming DPO (56.1\%) by 30.6 pp, GRPO (36.1\%) by 50.6 pp, and PPO (20.4\%) by 66.3 pp.

    \item \textbf{Correct directional adaptation}: On TruthfulQA, CITA improves by \textbf{+0.054} (adjusting uncertainty appropriately), while DPO shows minimal change (\textbf{+0.001}). This validates CITA's ability to internalize calibration instructions.

    \item \textbf{Mandatory KL is essential}: The KL term---optional in DPO---becomes mandatory in CITA to prevent mode collapse and enable stable instruction-conditioned behavioral switching.

    \item \textbf{Instruction-alignment $\neq$ instruction-following}: Through the ECLIPTICA benchmark, we demonstrate that CITA achieves \textit{deeper} instruction-alignment (internalized policies) rather than \textit{superficial} instruction-following (surface compliance).
\end{enumerate}

\vspace{-2mm}
\subsection{Implications}
\label{sec:implications}

CITA enables a new paradigm of \textbf{interactive and adaptive AI systems} where:

\begin{itemize}[leftmargin=*,itemsep=2pt]
    \item \textbf{Policy updates via natural language}: Alignment changes can be expressed through instructions without model retraining.

    \item \textbf{Context-aware responses}: Different users, applications, or deployment contexts can receive appropriately calibrated behaviors.

    \item \textbf{Dynamic safety policies}: Safety constraints can be adjusted at inference time based on deployment requirements (e.g., stricter for child-facing apps, more permissive for research contexts).

    \item \textbf{Reduced alignment tax}: No need for multiple specialized models for different alignment policies---a single CITA model can switch behaviors on demand.
\end{itemize}

\textbf{Outlook.} CITA opens several research directions: (1) scaling to larger models~\cite{touvron2023llama,jiang2023mistral,team2023gemini,chowdhery2022palm} and more instruction types, (2) theoretical analysis of instruction-conditioned learning dynamics~\cite{azar2023general}, (3) applications to multi-stakeholder alignment where different parties have different policy requirements~\cite{sun2024trustllm}, and (4) integration with constitutional AI methods~\cite{bai2022constitutional} for hierarchical instruction handling.

\vspace{-2mm}
\subsection{Reproducibility}
\label{sec:reproducibility}

All code, trained models, and the ECLIPTICA benchmark are publicly available:

\begin{itemize}[leftmargin=*,itemsep=1pt]
    \item \textbf{Code}: \url{https://anonymous.4open.science/r/CITA_Anonymous-AC02}
    \item \textbf{ECLIPTICA benchmark}: \url{https://huggingface.co/datasets/anonymousML123/ISD-Instruction-Switch-Dataset}
\end{itemize}







\textbf{Benefits of CITA}
\label{sec:benefits}

\begin{enumerate}[leftmargin=*,itemsep=2pt]
    \item \textbf{Generalization}: Responds appropriately to novel instruction-prompt combinations not seen during training, similar to zero-shot generalization in instruction tuning~\cite{wei2022finetuned,chung2022scaling}.

    \item \textbf{Robustness}: Resists jailbreak attempts~\cite{zou2023universal,wei2023jailbroken} by conditioning behavior on explicit alignment instructions.

    \item \textbf{Calibration}: Avoids over-refusal~\cite{bianchi2023safetytuned} by adjusting safety level based on instruction context.

    \item \textbf{Behavioral Switching}: Dynamically adapts response style (e.g., concise vs. detailed, formal vs. casual) based on instructions, enabling multi-tenant deployment~\cite{wang2024survey}.
\end{enumerate}

































% =============================================================================
% NOTE: Table "CITA vs. Prior Methods" moved to Section 3 (4_cita_framework.tex)
% See Table~\ref{tab:feature_comparison} in Section~\ref{sec:cita_framework}
% =============================================================================
